% !TeX root = ../drafts/b-polynomial-commitment-scheme.tex
\section{The polynomial commitment scheme: WHIR \slash Ligerito}\label{annex:pcs}

A self-contained description of the WHIR\slash Ligerito polynomial commitment scheme \cite{ACFY25,NA25} over binary fields, with a round-by-round soundness analysis in the Johnson list-decoding regime. The two schemes coincide when WHIR is instantiated with interleaved codes and Ligerito with the Reed-Solomon codes of Annex~\ref{app:ntt}, which gives the novel polynomial basis and additive NTT used for encoding. Following Flock~\cite{Flock26}, we do not use out-of-domain (OOD) sampling at commitment, so our PCS is only list binding. The two fields are those of \S\ref{sec:vm}: we commit to a $\K$-valued multilinear, and opening claims are in $\E$.

\paragraph{Symbols.}

\begin{center}\small
\begin{tabular}{@{}lll@{}}
\hline
here & in \cite{ACFY25,NA25,BCHKS25} & meaning\\
\hline
$\nv_i$        & $m_i$        & variables of the multilinear at level $i$\\
$\nlev$        & $r$          & number of levels\\
$\ell_i$       & $\ell_i$, $\nu$ & fold factor, and the interleaving exponent\\
$\rate_i$, $\rho_{-,i}$ & $\rho_i$, $(k_i-1)/n_i$ & code rate and reduced degree ratio\\
$\rexp_i$      & $R_i$        & rate exponent, $\rate_i=2^{-\rexp_i}$\\
$\rad_i$       & $\gamma_i$   & proximity radius\\
$\slack_i$     & $\eta_i$     & Johnson slack\\
$\mcapar$      & $m$          & MCA interpolation multiplicity\\
$\dmin$        & $\delta$     & relative minimum distance\\
$\lambda$      & $\lambda$, $\alpha$ & batching challenge, as in \S\ref{sec:stacking}\\
$\fc_j$        & $s_j$        & sumcheck fold challenge\\
$\rc{j}$       & $\sigma_j$   & running sumcheck claim after round $j$\\
$J_i$          & $T_i$        & number of claims, as in \S\ref{sec:e2e-unrolled}\\
$\nq_i$        & $t_i$        & query count\\
$\qset_i$      & $J_i$        & the sampled columns\\
$\blen$        & $n$          & code block length\\
$\dom$         & $S$          & evaluation domain\\
$\orc^{(i)}$   & $C^{(i)}$    & oracle, an interleaved word ($C_j$ is a constraint)\\
$\nrows$       & (none)       & rows of $\orc^{(0)}$ the prover commits, $\nrows \le 2^{\ell_0}$\\
$\nstack$      & (none)       & field elements placed in the stack, as in \S\ref{sec:stacking}\\
$\lst_{\rad}$  & $\lambda_\gamma$ & list of codewords within $\rad$\\
$W$, $c$       & $w$, $v$     & weight and value of a claim, as in \S\ref{sec:stacking}\\
$\ffinal$      & $p$          & final plaintext multilinear\\
$\cood$        & $y$          & out-of-domain answer ($y$ generates $\E$ over $\K$)\\
$\mcanum$      & $a$          & numerator of the MCA error\\
$\vansub_i$, $\Wn{i}$ & $W_i$, $\widehat W_i$ & subspace vanishing polynomial, normalized\\
$\novel_j$     & $X_j$        & novel-basis polynomial ($X_i$ is a sumcheck variable)\\
$\ntmap_d$     & $q_d$        & linearized map of the NTT ($q$ is the witness stack)\\
$\bfarr$       & $B$          & butterfly array\\
$\mpoly$       & $P$          & message polynomial ($P_i$ is a committed column)\\
$e_c$          & $\beta_c$    & the fixed $\Ftwo$-basis, as in Annex \ref{annex:rs}\\
\hline
\end{tabular}
\end{center}

%=====================================================================
\subsection{Preliminaries}\label{sec:prelims}
%=====================================================================

\subsubsection{Codes, interleaving, and list decoding}\label{sec:codes}

We use the Reed-Solomon code $\RS[\E,\dom,\kdim]$ of Definition~\ref{def:rs-code}, with $|\dom|=\blen$, dimension $\kdim=2^\kappa$, rate $\rate=\kdim/\blen$, degree bound $D=\kdim-1$, and reduced degree ratio $\rho_-=D/\blen$. Messages are indexed by $\cube{\kappa}$; the linear encoder $\Enc$ of \S\ref{sec:binary-rs} maps them onto the code and preserves $\K$-valued words.

\begin{definition}[Interleaved words, lists, folds]\label{def:interleaved}
An \emph{interleaved word} is a matrix $\orc \in \E^{\cube{\ell} \times \dom}$; it is an \emph{interleaved codeword} if every row $\orc[u, \cdot]$ is a codeword $(u \in \cube{\ell})$. $\orc$ \emph{agrees} with an interleaved codeword $U$ on $\aset \subseteq \dom$ if $\orc[u,x] = U[u,x]$ for all $u \in \cube{\ell}$ and all $x \in \aset$; it is \emph{$\rad$-close} to $U$ if they agree on some $\aset$ with $|\aset| \ge (1-\rad)\blen$; and $\lst_\rad(\orc)$ is the set of interleaved codewords $\rad$-close to $\orc$. A single word \emph{agrees with the code} on $\aset$ if some codeword equals it there, and $\orc$ \emph{agrees with the interleaved code} on $\aset$ if some interleaved codeword equals its rows there. An interleaved codeword $U$ \emph{decodes} to the multilinear $g_U : \cube{\ell + \kappa} \to \E$ with $U[u, \cdot] = \Enc(g_U(u, \cdot))$; distinct $U$ decode to distinct $g_U$, hence distinct $\mle{g_U}$. For $\fc \in \E^{j}$, $j \le \ell$, the \emph{fold} $\orc_\fc \in \E^{\cube{\ell - j} \times \dom}$ is $\orc_\fc[u, x] := \sum_{\iota \in \cube{j}} \eq(\fc, \iota)\, \orc[(\iota, u), x]$ (for $j = \ell$, a single word in $\E^{\dom}$).
\end{definition}

\begin{fact}[Folding commutes with decoding]\label{fact:fold-decode}
If $U$ is an interleaved codeword then so is $U_\fc$, with $\mle{g_{U_\fc}} = \mle{g_U}(\fc, \cdot)$ (linearity of $\Enc$ plus Fact~\ref{fact:eq}).
\end{fact}

\begin{theorem}[Johnson bound]\label{thm:johnson-interleaved}
For every interleaved word $\orc$ of $\RS[\E,\dom,\kdim]$, and every integer agreement threshold $A$ satisfying $A^2>\blen(\kdim-1)$,
\[
 |\lst_{1-A/\blen}(\orc)| \le \left\lfloor\frac{\blen(A-\kdim+1)}{A^2-\blen(\kdim-1)}\right\rfloor.
\]
See \S\ref{app:johnson} for a self-contained proof. The bound is independent of the interleaving width because a column can be encoded injectively as one symbol over a field extension, preserving equality.
\end{theorem}

\subsubsection{PCS and round-by-round soundness}\label{sec:iopcs}

\begin{definition}[$L$-list binding, informal]\label{def:listbinding}
A \emph{polynomial commitment scheme} (PCS) is \emph{$L$-list binding} if every commitment transcript $\mathsf{cm}$ determines a set $\polyset_{\mathsf{cm}}$ of at most $L$ multilinear polynomials, not necessarily efficiently computable, such that for every unbounded prover and every collection of opened claims: if no member of $\polyset_{\mathsf{cm}}$ satisfies all the claims, the verifier rejects except with small probability. $L = 1$ is standard binding.
\end{definition}

\begin{definition}[RBR soundness, informal; cf.~\cite{CCHLRR19,ACFY25}]\label{def:rbr}
A public-coin protocol is \emph{round-by-round (RBR) sound with errors} $(\varepsilon_1, \varepsilon_2, \dots)$, one per verifier message, if there is an \emph{invariant} on partial transcripts (not necessarily efficiently computable) marking the positions where a cheating prover has already lost:
\begin{enumerate}[itemsep=1pt,label=(\roman*)]
\item on a false statement, the invariant holds before any message is sent;
\item the prover alone cannot escape it: it survives every prover message;
\item the $i$-th verifier message lets it escape with probability at most $\varepsilon_i$;
\item if it survives to the end of the interaction, the verifier rejects.
\end{enumerate}
\end{definition}

Interactive soundness error is then at most $\sum_i \varepsilon_i$. After Fiat--Shamir, it is $\max_i \varepsilon_i$ per random-oracle query~\cite{CCHLRR19,BCS16}.

%=====================================================================
\subsection{Reed--Solomon codes over binary fields}\label{sec:binary-rs}
%=====================================================================

The encoder $\Enc_{\kappa,\rexp}$ of Definition~\ref{def:enc} evaluates $2^\kappa$ novel-basis coefficients on an additive domain of size $\blen=2^{\kappa+\rexp}$. Here $\rexp\ge1$ and $\kappa+\rexp\le64$. The protocol uses two properties, proved in Annex~\ref{app:ntt}:
\begin{itemize}[itemsep=2pt]
\item \emph{fast encoding, for the prover:} $\Enc_{\kappa,\rexp}$ is computed by the additive NTT with $\tfrac12 \kappa \blen$ multiplications and fewer than $(\kappa + 1)\, \blen$ additions (Proposition~\ref{prop:ntt});
\item \emph{MLE-friendly columns, for the verifier:} for every $x \in \dom$, the weight $W_x : \cube{\kappa} \to \K$, defined by $\Enc_{\kappa,\rexp}(g)[x] = \inner{W_x}{g}$ (the relation being valid for arbitrary $g$), admits an MLE efficiently evaluable anywhere in $\E^{\kappa}$ in $O(\kappa)$ field operations (Lemma~\ref{lem:colweight}).
\end{itemize}

%=====================================================================
\subsection{The protocol}\label{sec:protocol}
%=====================================================================

\begin{protocol}[Whir/Ligerito]\label{fig:protocol}
\ \par\smallskip
\emph{Parameters.} The scheme commits to $f : \cube{\nv} \to \K$ across $\nlev \ge 1$ \emph{levels} $i = 0, \dots, \nlev-1$ plus a plaintext final level $\nlev$: variable counts $\nv = \nv_0 > \nv_1 > \dots > \nv_\nlev \ge 0$ with fold factors $\ell_i := \nv_i - \nv_{i+1} \ge 1$ and $\nv_\nlev$ small (e.g.\ $\nv_\nlev \le 5$); per level, a rate $\rate_i = 2^{-\rexp_i}$, $\rexp_i \ge 1$, giving the encoder $\Enc_i := \Enc_{\kappa_i, \rexp_i}$ of \S\ref{sec:binary-rs} with $\kappa_i := \nv_{i+1}$, i.e.\ the code $\mathcal{C}_i := \RS[\E, \dom_i, 2^{\kappa_i}]$ over its domain $\dom_i \subseteq \K$ of size $\blen_i = 2^{\kappa_i + \rexp_i}$ (requires $\kappa_i + \rexp_i \le 64$); and a query count $\nq_i$. Its degree bound is $D_i:=2^{\kappa_i}-1$ and its reduced degree ratio is $\rho_{-,i}:=D_i/\blen_i$, distinct from the code rate $\rate_i=2^{\kappa_i}/\blen_i$. The integer agreement thresholds $A_i>\sqrt{D_i\blen_i}$ appear only in the analysis.

\smallskip
\textbf{Commit}$(f : \cube{\nv} \to \K)$: $\mathbf{P}$ sends the oracle $\orc^{(0)}[u, \cdot] := \Enc_0(f(u, \cdot))$, $u \in \cube{\ell_0}$, whose entries are $\K$ elements. No OOD sample is taken.

\smallskip
\textbf{Open}$\bigl(\{(W_t, c_t)\}_{t = 1, \dots, J_0}\bigr)$, $J_0 \ge 1$ weighted claims about $f$: set $f^{(0)} := f$. For $i = 0, 1, \dots, \nlev - 1$:
\begin{enumerate}[itemsep=0.5pt,topsep=2pt]
\item \emph{(batch)} $\mathbf{V}$ sends $\lambda \leftarrow \E$; the \emph{batched claim} is $(W, \rc{0}) := \bigl(\sum_{t = 1}^{J_i} \lambda^{t-1} W_t,\; \sum_{t} \lambda^{t-1} c_t\bigr)$.
\item \emph{(sumcheck~\cite{LFKN92})} For $j = 1, \dots, \ell_i$: $\mathbf{P}$ sends $h_j \in \E[X]$ with $\deg h_j \le 2$, honestly $h_j(X) = \sum_{x} \mle{W}(\fc^{<j}, X, x)\, \mle{f^{(i)}}(\fc^{<j}, X, x)$ where $\fc^{<j} := (\fc_1, \dots, \fc_{j-1})$; $\mathbf{V}$ checks $h_j(0) + h_j(1) \qeq \rc{j-1}$, samples $\fc_j \leftarrow \E$ and sends it to the prover; both set $\rc{j} := h_j(\fc_j)$.\\
Set $\fc := (\fc_1, \dots, \fc_{\ell_i})$, $c' := \rc{\ell_i}$, and, for $x \in \cube{\nv_{i+1}}$: \ $f^{(i+1)}(x) := \mle{f^{(i)}}(\fc, x)$, \ $W'(x) := \mle{W}(\fc, x)$. The \emph{residual claim} $(W', c')$ is a weighted claim about the folded multilinear $f^{(i+1)} : \cube{\nv_{i+1}} \to \E$.
\item If $i < \nlev - 1$:
\begin{enumerate}[itemsep=0.5pt,topsep=1pt,label=(\alph*)]
\item \emph{(commit)} $\mathbf{P}$ sends the oracle $\orc^{(i+1)}[u, \cdot] := \Enc_{i+1}(f^{(i+1)}(u, \cdot))$, $u \in \cube{\ell_{i+1}}$.
\item \emph{(ood)} $\mathbf{V}$ sends $z \leftarrow \E^{\nv_{i+1}}$; $\mathbf{P}$ sends $\cood \in \E$ (honestly $\mle{f^{(i+1)}}(z)$).
\item \emph{(query)} $\mathbf{V}$ sends $\nq_i$ i.i.d.\ columns $\qset_i = (x_q)_{q = 1, \dots, \nq_i} \leftarrow \dom_i$; it reads $\orc^{(i)}[\cdot, x_q]$ and computes $c_q := \sum_{u} \eq(\fc, u)\, \orc^{(i)}[u, x_q] = \orc^{(i)}_\fc[x_q]$; honestly $c_q = \Enc_i(f^{(i+1)})[x_q]$, by Fact~\ref{fact:fold-decode}.
\item Claims for level $i+1$ (so $J_{i+1} := \nq_i + 2$): the residual claim $(W', c')$, the OOD claim $(\eq(z, \cdot), \cood)$, and the $\nq_i$ \emph{consistency claims} $(W^{(i)}_{x_q}, c_q)$, where $W^{(i)}_x$ are the column weights of $\Enc_i$ (\S\ref{sec:binary-rs}): the $q$-th asserts $\Enc_i(f^{(i+1)})[x_q] = c_q$.
\end{enumerate}
\item If $i = \nlev - 1$: \emph{(final)} $\mathbf{P}$ sends the multilinear $\ffinal : \cube{\nv_\nlev} \to \E$; $\mathbf{V}$ sends $\nq_{\nlev-1}$ i.i.d.\ columns $\qset_{\nlev-1} \leftarrow \dom_{\nlev-1}$ and computes $c_q$ from $\orc^{(\nlev-1)}$ as in 3(c). Consistency is again a weighted claim, $\Enc_{\nlev-1}(\ffinal)[x_q] = \inner{W^{(\nlev-1)}_{x_q}}{\ffinal}$ (Lemma~\ref{lem:colweight}), so the level closes like the others: $\mathbf{V}$ samples and sends $\lambda \leftarrow \E$, batching the $\nq_{\nlev-1}$ consistency claims with the residual claim $(W', c')$ into one claim about $\ffinal$; $\nv_\nlev$ further sumcheck rounds discharge it, and $\mathbf{V}$ closes by evaluating $\mle{\ffinal}$ at their point from the table it holds.
\end{enumerate}
$\mathbf{V}$ accepts iff every check above passed.
\end{protocol}

All weights in the protocol are MLE-friendly and multilinear, so every sumcheck summand $\mle{W} \cdot \mle{f^{(i)}}$ has individual degree at most $2$. At level $0$ that summand pairs an $\E$-valued weight with the $\K$-valued $f$; every later level is $\E$ throughout.

Completeness is perfect: the honest prover passes every check ($h_j(0) + h_j(1) = \rc{j-1}$ by definition of $h_j$, $c_q = \Enc_i(f^{(i+1)})[x_q]$ by Fact~\ref{fact:fold-decode}, the OOD and residual claims by construction, the final checks as identities).

%=====================================================================
\subsection{Round-by-round soundness}\label{sec:soundness}
%=====================================================================

Fix a commitment $\mathsf{cm} = \orc^{(0)}$ and input claims $\{(W_t, c_t)\}_t$, and define the relation
\[
\relopen \;:=\; \Bigl\{ \bigl(\orc^{(0)}, \{(W_t, c_t)\}_t\bigr) \;:\; \exists\, U \in \lst_{\rad_0}(\orc^{(0)}) \text{ with } \inner{W_t}{g_U} = c_t \text{ for all } t \Bigr\}.
\]
The commitment binds the prover to the list $\polyset_{\mathsf{cm}} := \{\mle{g_U} : U \in \lst_{\rad_0}(\orc^{(0)})\}$. We set $\rad_i:=1-A_i/\blen_i$ and write $L_i$ for the exact integer list bound in Theorem~\ref{thm:mca-johnson}; in particular, $|\polyset_{\mathsf{cm}}|\le L_0$.

\begin{theorem}[RBR soundness]\label{thm:rbr}
Suppose that at every level $i$, the integer threshold $A_i$ and interpolation certificate satisfy Theorem~\ref{thm:mca-johnson}. Let $L_i$ be its list bound, let $E_i$ be its exceptional count, and set $\rad_i:=1-A_i/\blen_i$. Then Protocol~\ref{fig:protocol} is an $L_0$-list binding PCS, and its opening phase for $\relopen$ has the following RBR soundness error at each verifier message:
\[
\begin{array}{llcl}
\textnormal{batching, level } i: & \varepsilon^{\mathrm{batch}}_i &=& \dfrac{(J_i - 1)\, L_i}{|\E|} , \\[2ex]
\textnormal{fold challenge } \fc_j, \textnormal{ level } i: & \varepsilon^{\mathrm{fold}}_{i,j} &=& \dfrac{E_i+2 L_i}{|\E|} , \\[2ex]
\textnormal{OOD challenge entering level } i \ge 1: & \varepsilon^{\mathrm{ood}}_i &=& \dbinom{L_i}{2} \cdot \dfrac{\nv_i}{|\E|} , \\[2ex]
\textnormal{query message } \qset_i: & \varepsilon^{\mathrm{query}}_i &=& \left(\dfrac{A_i-1}{\blen_i}\right)^{\nq_i} , \\[2ex]
\textnormal{final batching challenge } \lambda: & \varepsilon^{\mathrm{fin}} &=& \dfrac{\nq_{\nlev-1}}{|\E|} , \\[2ex]
\textnormal{each of the } \nv_\nlev \textnormal{ closing rounds}: & \varepsilon^{\mathrm{tail}} &=& \dfrac{2}{|\E|} .
\end{array}
\]
Here $J_0$ is the number of input claims and $J_i = \nq_{i-1} + 2$ for $i \ge 1$. The \emph{RBR error}, which governs the Fiat--Shamir compilation (\S\ref{sec:iopcs}), is the maximum of these entries over all rounds.
\end{theorem}

The remainder of this section proves the theorem: \S\ref{sec:tools} states the coding-theoretic tools, \S\ref{pcs:proof} the invariant and the round-by-round argument.

\subsubsection{Proximity tools}\label{sec:tools}

The engine is \emph{mutual correlated agreement} (MCA)~\cite{ACFY25}. For a two-row word $\orc$, recall $\orc_\fc = (1-\fc)\orc[0,\cdot] + \fc\,\orc[1,\cdot]$.\footnote{\cite{BCHKS25} state the results for the affine line $\orc[0,\cdot] + \fc\,\orc[1,\cdot]$; in characteristic $2$, $\orc_\fc = \orc[0,\cdot] + \fc(\orc[0,\cdot] + \orc[1,\cdot])$, and this invertible change of pair maps codeword pairs to codeword pairs, so the formulations are equivalent.}

\begin{theorem}[MCA in the unique-decoding regime; adapted from {\cite[Cor.~1.4]{BCHKS25}}]\label{thm:mca-udr}
Let $\RS[\E,\dom,\kdim]$ have block length $\blen = |\dom|$ and relative minimum distance $\dmin := 1 - (\kdim-1)/\blen$. Suppose
\[
\dmin \ge 3\sqrt{\frac{2}{\blen}}
\qquad\text{and}\qquad
\rad \in \left[\frac{\dmin}{3},\,\frac{\dmin}{2} - \frac{3}{\dmin \blen}\right].
\]
Then, for every two-row word $\orc$,
\[
\Pr_{\fc \leftarrow \E}\!\left[\; \exists \aset \subseteq \dom,\ |\aset| \ge (1-\rad) \blen :\ \begin{array}{l} \orc_\fc \text{ agrees with the code on } \aset, \text{ and} \\ \orc \text{ does not agree with the interleaved code on } \aset \end{array} \right]
\;\le\; \frac{\rad \blen + 1}{|\E|}.
\]
\end{theorem}

We record Theorem~\ref{thm:mca-udr} only for comparison, we don't need it in the List-Decoding Regime, we instead need:

\begin{theorem}[MCA up to the Johnson bound; {\cite{DKT26}}]\label{thm:mca-johnson}
Let $\RS[\E,\dom,\kdim]$ have block length $\blen=|\dom|$, degree bound $D:=\kdim-1$, and reduced degree ratio $\rho_-:=D/\blen$, where $1\le D\le\blen-2$. Fix an integer threshold $D<A\le\blen$ with $A^2>\blen D$. Choose integers $m\ge1$, $1\le B\le\lfloor(mA-1)/D\rfloor$, and $H\ge B$. Set $u:=\min\{B,m-1\}$ and
\[
 N:=\sum_{b=0}^{B}(mA-Db),\qquad W:=\sum_{b=0}^{B}b(mA-Db),\qquad R:=\sum_{b=0}^{u}(m-b),\qquad T:=\sum_{b=0}^{u}b(m-b).
\]
Suppose $N>\blen R$ and $(H+1)N-W>\blen((H+1)R-T)$. Define
\[
 \Psi_D(B):=1+(2D-1)(2B-1)+2(B-2D-1)_+
\]
and
\[
 E:=(2B-1)H+\frac{\blen-D}{A-D}\bigl(B+H\Psi_D(B)\bigr)+(\blen-D-1)B,
 \qquad
 L:=\min\left\{B,\left\lfloor\frac{\blen(A-D)}{A^2-\blen D}\right\rfloor\right\}.
\]
Every received word has at most $L$ degree-at-most-$D$ codewords agreeing on at least $A$ positions. For every two-row word $\orc$,
\[
\Pr_{\fc \leftarrow \E}\!\left[\; \exists \aset \subseteq \dom,\ |\aset| \ge A :\ \begin{array}{l} \orc_\fc \text{ agrees with the code on } \aset, \text{ and} \\ \orc \text{ does not agree with the interleaved code on } \aset \end{array} \right] \;\le\; \frac{E}{|\E|}.
\]
The conclusion preserves the full common agreement set and holds over every field extension. A closed sufficient certificate sets $\slack:=A/\blen-\sqrt{\rho_-}>0$ and
\[
 m=\max\left\{\left\lceil\frac{\sqrt{\rho_-}}{2\slack}\right\rceil,3\right\},\qquad t=m+\frac12,\qquad B=\left\lceil\frac{t}{\sqrt{\rho_-}}\right\rceil-1,\qquad H=\left\lceil\frac{t^2}{3\rho_-}\right\rceil-1.
\]
\end{theorem}

The interpolation threshold is the classical Johnson threshold; the gain is in the exceptional count. At fixed reduced rate and positive Johnson slack, Theorem~\ref{thm:mca-johnson} gives $E=O(\blen/\slack^3)$. The bound of BCHKS~\cite[Thm.~4.6]{BCHKS25} uses $m_{\mathrm B}=\max\{\lceil\sqrt{\rho_-}/\slack\rceil,3\}$, $t_{\mathrm B}=m_{\mathrm B}+1/2$, and
\[
 E_{\mathrm{BCHKS}}=\frac{2t_{\mathrm B}^5+3t_{\mathrm B}(1-A/\blen)\rho_-}{3\rho_-^{3/2}}\blen+\frac{t_{\mathrm B}}{\sqrt{\rho_-}},
\]
which scales as $O(\blen/\slack^5)$. The implementation checks the displayed integer certificate directly, including the one finite certificate used instead of the closed choice.

\begin{lemma}[Folding preserves lists]\label{lem:fold-list}
Let $\RS[\E,\dom,\kdim]$ have full agreement-set line MCA at threshold $A$ with exceptional count $E$, as in Theorem~\ref{thm:mca-johnson}. For every interleaved word $\orc \in \E^{\cube{\ell} \times \dom}$, $\ell \ge 1$:
\[
\Pr_{\fc \leftarrow \E}\Bigl[\, \lst_{1-A/\blen}(\orc_\fc) \neq \{ U_\fc : U \in \lst_{1-A/\blen}(\orc) \} \,\Bigr] \;\le\; \frac{E}{|\E|} .
\]
The bound is independent of $\ell$.
\end{lemma}

\begin{proof}
$\supseteq$ holds always: folding an interleaved codeword that agrees on a common set of at least $A$ columns preserves those agreements. For $\subseteq$, regard all $2^{\ell-1}$ row pairs as one interleaved received line. The width-free interleaving transfer of~\cite{DKT26} applies the scalar full agreement-set MCA theorem once to this entire line, with exceptional count $E$. Outside that event, every codeword tuple near the folded word decomposes into two codeword tuples near the two coefficient arrays with the same full common agreement set. These tuples form an interleaved codeword $U$ near $\orc$, and their affine combination is the original folded candidate. Hence the two lists are equal.
\end{proof}

Applying the lemma successively to a height-$h$ multilinear fold gives ordinary interactive error at most $hE/|\E|$ by a union bound over its independent challenges. For RBR soundness, each current fold challenge is conditioned on the preceding transcript and pays only $E/|\E|$.

\begin{lemma}[OOD separation]\label{lem:ood}
Let $\orc$ be an interleaved word with $|\lst_\rad(\orc)| \le L$, decoding to $\nv$-variable multilinears. Then $\Pr_{z \leftarrow \E^{\nv}}[\exists \text{ distinct } U, U' \in \lst_\rad(\orc) : \mle{g_U}(z) = \mle{g_{U'}}(z)] \le \binom{L}{2} \nv / |\E|$.
\end{lemma}

\begin{proof}
Distinct interleaved codewords have distinct multilinear extensions; apply Lemma~\ref{lem:sz} to each difference and union bound.
\end{proof}

\subsubsection{Proof of the main theorem}\label{pcs:proof}

At level $i$, abbreviate $\fc^{\le j} := (\fc_1, \dots, \fc_j)$ and $\fc^{<j} := \fc^{\le j-1}$; thus $\fc^{\le 0}$ is empty and $\orc^{(i)}_{\fc^{\le 0}} = \orc^{(i)}$. For $0 \le j \le \ell_i$, the \emph{round-$j$ claim} is the weighted claim
\[
\bigl(\, \mle{W}(\fc^{\le j}, \cdot),\; \rc{j} \,\bigr) \qquad \text{about multilinears } g : \cube{\nv_i - j} \to \E,
\]
where $W$, $\rc{0}$, and $\rc{j} = h_j(\fc_j)$ are read off the transcript as in Protocol~\ref{fig:protocol}; an interleaved codeword satisfies a claim if its decoded multilinear does. The round-$0$ claim is the batched claim, and the round-$\ell_i$ claim is the residual claim $(W', c')$.

We now define the invariant required by Definition~\ref{def:rbr}, a condition on the partial transcripts of the opening. In one phrase: \emph{every codeword close to the current fold violates the current claim}. Precisely, the list below gives the condition just after each verifier message; a transcript ending with prover messages inherits the condition of its last verifier message. Every condition also includes one implicit disjunct, the \emph{escape clause}: \emph{``some verifier check has already failed''} (such transcripts are rejected no matter what follows).
\begin{itemize}[itemsep=1pt]
\item \emph{Start of level $i$} (before any message for $i = 0$, after the query message $J_{i-1}$ otherwise): every $U \in \lst_{\rad_i}(\orc^{(i)})$ violates at least one of the level's $J_i$ claims (the input claims for $i = 0$, the claims of step 3(d) of Protocol~\ref{fig:protocol} otherwise).
\item \emph{After the batching challenge $\lambda$ of level $i$, and after each of its fold challenges $\fc_j$:} every $U \in \lst_{\rad_i}(\orc^{(i)}_{\fc^{\le j}})$ violates the round-$j$ claim ($\lambda$ being the case $j = 0$).
\item \emph{After the OOD challenge $z$ closing level $i$ ($i < \nlev-1$):} the round-$\ell_i$ instance above holds, and the multilinear extensions decoded from $\lst_{\rad_{i+1}}(\orc^{(i+1)})$ are pairwise distinct at $z$.
\item \emph{After the final query message $\qset_{\nlev-1}$:} $\ffinal$ violates one of the level's $\nq_{\nlev-1}+1$ claims.
\item \emph{After the final $\lambda$, and after each of the $\nv_\nlev$ rounds that follow (complete transcript):} $\ffinal$ violates the running claim. The terminal one $\mathbf{V}$ checks against $\mle{\ffinal}$ itself, so it rejects.
\end{itemize}
Conditions (i), (ii), (iv) of Definition~\ref{def:rbr} hold by inspection: on a false statement, the first item is exactly the failure of $\relopen$; a prover message changes nothing an instance depends on, except possibly failing a check, which only switches the escape clause on; and the last item is precisely ``some check failed'', so such a complete transcript is rejected. The separation clause of the OOD instance is there because the query message consumes it: it pins the new oracle's list to at most one candidate consistent with the prover's answer. It remains to bound the escape probabilities (iii); in the proof below we assume all checks so far pass, since otherwise the escape clause already holds.

\begin{proof}[Proof of Theorem~\ref{thm:rbr}]
If the statement lies outside $\relopen$, then before any message is sent every $U \in \lst_{\rad_0}(\orc^{(0)})$ violates at least one input claim, so the invariant holds. We follow the protocol and bound, for each verifier message, the probability that the invariant fails to carry over. Every list below consists of interleaved words of $\mathcal{C}_i$ within radius $\rad_i$, hence has size at most $L_i$ by Theorems~\ref{thm:mca-johnson} and~\ref{thm:johnson-interleaved}.

\medskip\noindent\emph{Batching challenge $\lambda$.}
Assume that every $U \in \lst_{\rad_i}(\orc^{(i)})$ violates at least one of the $J_i$ claims $(W_t, c_t)$. Fix such a $U$. Its discrepancy against the batched claim is $\sum_{t} \lambda^{t - 1} d_t$ with $d_t := \inner{W_t}{g_U} - c_t$, a polynomial in $\lambda$ of degree less than $J_i$, nonzero because the $d_t$ are not all zero. By Lemma~\ref{lem:sz} it vanishes with probability at most $(J_i - 1)/|\E|$, so a union bound over the at most $L_i$ list members leaves every $U$ violating the round-$0$ claim except with probability $(J_i - 1) L_i / |\E|$.

\medskip\noindent\emph{Fold challenge $\fc_j$.}
Assume that every $U \in \lst := \lst_{\rad_i}(\orc^{(i)}_{\fc^{<j}})$ violates the round-$(j{-}1)$ claim. Since all checks so far pass (otherwise the verifier trivially rejects), $h_j(0) + h_j(1) = \rc{j-1}$. For $U \in \lst$, let
\[
H_U(X) \;:=\; \sum_{x \in \cube{\nv_i - j}} \mle{W}(\fc^{<j}, X, x)\, \mle{g_U}(X, x)
\]
be the honest round polynomial for $U$; it has degree at most $2$, since $\mle{W}$ and $\mle{g_U}$ are multilinear. Summing over $X \in \{0, 1\}$ turns the multilinear evaluations back into values, so
\[
H_U(0) + H_U(1) \;=\; \bigl\langle\, \mle{W}(\fc^{<j}, \cdot),\; g_U \,\bigr\rangle \;\neq\; \rc{j-1} \;=\; h_j(0) + h_j(1),
\]
the inequality being the assumed violation; hence $h_j \neq H_U$ as polynomials. Two bad events over $\fc_j$:
\begin{itemize}[itemsep=1pt]
\item $B_1$: $h_j(\fc_j) = H_U(\fc_j)$ for some $U \in \lst$. For each $U$, two distinct polynomials of degree at most $2$ agree on at most $2$ points, so $\Pr[B_1] \le 2 L_i / |\E|$.
\item $B_2$: $\lst_{\rad_i}(\orc^{(i)}_{\fc^{\le j}}) \neq \{ U_{\fc_j} : U \in \lst \}$, i.e.\ folding does not map the old list onto the new one. Lemma~\ref{lem:fold-list} applies to all remaining rows as one interleaved word and gives $\Pr[B_2] \le E_i/|\E|$, independently of their number.
\end{itemize}
If neither occurs, every $V \in \lst_{\rad_i}(\orc^{(i)}_{\fc^{\le j}})$ is $V = U_{\fc_j}$ for some $U \in \lst$, with $\mle{g_V} = \mle{g_U}(\fc_j, \cdot)$ by Fact~\ref{fact:fold-decode}, so
\[
\bigl\langle\, \mle{W}(\fc^{\le j}, \cdot),\; g_V \,\bigr\rangle
\;=\; \sum_{x \in \cube{\nv_i - j}} \mle{W}(\fc^{<j}, \fc_j, x)\, \mle{g_U}(\fc_j, x)
\;=\; H_U(\fc_j) \;\neq\; h_j(\fc_j) \;=\; \rc{j} ,
\]
using $\neg B_1$ for the inequality: every $V$ violates the round-$j$ claim. The combined fold error is $(2 L_i+E_i)/|\E|$.

\medskip\noindent\emph{The level interface.}
After the last fold challenge of level $i < \nlev - 1$, we know: every codeword within $\rad_i$ of the single row $c := \orc^{(i)}_\fc \in \E^{\dom_i}$ violates the residual claim $(W', c')$. The prover now sends the new oracle $\orc^{(i+1)}$.

\medskip\noindent\emph{OOD challenge $z$.}
After $z$, two properties must hold: the one above, which does not involve $z$; and pairwise distinctness at $z$ of the multilinear extensions decoded from $\lst_{\rad_{i+1}}(\orc^{(i+1)})$, a list already determined by the prover's message. The latter fails with probability at most $\binom{L_{i+1}}{2} \nv_{i+1} / |\E|$ over $z$, by Lemma~\ref{lem:ood}.

\medskip\noindent\emph{Query message $\qset_i$.}
Assume both properties of the previous step; the prover has since sent its answer $\cood$. The values $\mle{g_U}(z)$ are pairwise distinct across $\lst_{\rad_{i+1}}(\orc^{(i+1)})$, so at most one member $U^\ast$ can satisfy $\mle{g_{U^\ast}}(z) = \cood$: every other member violates the OOD claim $(\eq(z, \cdot),\, \cood)$, whatever $\qset_i$ is. If no such $U^\ast$ exists, we are done. Otherwise write $g^\ast := g_{U^\ast}$, a multilinear on $\cube{\nv_{i+1}}$, and distinguish two cases.
\begin{itemize}[itemsep=1pt]
\item \emph{$\Enc_i(g^\ast)$ agrees with $c$ on at least $(1 - \rad_i) \blen_i$ coordinates.} Then the codeword $\Enc_i(g^\ast)$ lies in $\lst_{\rad_i}(c)$, hence violates the residual claim by assumption: $\inner{W'}{g^\ast} \neq c'$, whatever $\qset_i$ is.
\item \emph{Otherwise.} Because $\Enc_i(g^\ast)$ is outside the threshold-$A_i$ list, it agrees with $c$ on at most $A_i-1$ columns. All $\nq_i$ independent uniform queries land in this agreement set with probability at most $((A_i-1)/\blen_i)^{\nq_i}$; any disagreeing query $x_q$ makes $U^\ast$ violate the consistency claim $(W^{(i)}_{x_q},\, c_q)$, since $c_q = c[x_q]$.
\end{itemize}
Except with probability $((A_i-1)/\blen_i)^{\nq_i}$, every member of $\lst_{\rad_{i+1}}(\orc^{(i+1)})$ thus violates at least one of the $J_{i+1}$ claims, which is exactly the assumption of the batching step at level $i + 1$.

\medskip\noindent\emph{Final level.}
At level $\nlev - 1$ the fold challenges end as above: every codeword within $\rad_{\nlev-1}$ of $c := \orc^{(\nlev-1)}_\fc$ violates the residual claim. The prover sends $\ffinal$. If $\Enc_{\nlev-1}(\ffinal)$ agrees with $c$ on at least $A_{\nlev-1}$ coordinates then $\Enc_{\nlev-1}(\ffinal) \in \lst_{\rad_{\nlev-1}}(c)$, so $\ffinal$ violates the residual claim; otherwise, except with probability $((A_{\nlev-1}-1)/\blen_{\nlev-1})^{\nq_{\nlev-1}}$, some query disagrees and $\ffinal$ violates that consistency claim. Either way it violates one of the $\nq_{\nlev-1}+1$ batched claims, hence the batch itself except with probability $\nq_{\nlev-1}/|\E|$ (Lemma~\ref{lem:sz}), with no union over a list: $\ffinal$ is transmitted, not one of the codewords near an oracle. The $\nv_\nlev$ rounds that follow are a plain sumcheck (Fact~\ref{fact:sumcheck}) on a polynomial that $\mathbf{V}$ holds, so the terminal check against $\mle{\ffinal}$ fails.

\end{proof}


%=====================================================================
\subsection{Optimizations}\label{sec:pcs-opt}
%=====================================================================

The committed $f$ is the stack of \S\ref{sec:stacking}: $\nstack$ field elements in $\K$, then zeros up to $2^{\nv}$. Let the row index be the stack's most significant $\ell_0$ coordinates: $f(u, x) := q(x, u)$ for $u \in \cube{\ell_0}$ and $x \in \cube{\nv - \ell_0}$, \S\ref{app:ntt}. Row $u$ is then the block of $2^{\nv - \ell_0}$ consecutive elements of $q$ at offset $\langle u \rangle 2^{\nv - \ell_0}$, the padding is whole rows, and only $\nrows=\lceil\nstack/2^{\nv-\ell_0}\rceil$ of the $2^{\ell_0}$ rows hold data. So we get the following optimizations at commitment:
\begin{itemize}[itemsep=1pt,topsep=2pt]
\item the FFT-encoder only runs on $\nrows$ rows (and not the full $2^{\ell_0}$), since FFT-ing zeros gives zeros;
\item a Merkle leaf only contains $\nrows$ elements of $\K$ (proof size optimization), we keep the convention of hashing the full $2^{\ell_0}$ values for simplicity, but we hash them in reverse order, starting from the padding zeros;
\item the prover can thus preprocess the internal state of BLAKE2s after hashing those zeros, and then hash the useful part of each leaf starting from this state.
\end{itemize}

Conclusion: at commitment, both the FFT and the Merkle leaf hashing cost are proportional to $\nstack$ (not its next power of two). Unfortunately we still pay for the padding inside the $\nstack$ field elements, because some instruction counts are not powers of two. This could be prevented using the Jagged PCS \cite{jagged_pcs}.


%=====================================================================
\subsection{The Johnson bound}\label{app:johnson}
%=====================================================================

We prove Theorem~\ref{thm:johnson-interleaved} in its general form (arbitrary code).

\begin{theorem}[Johnson bound, agreement form]\label{thm:johnson-general}
Let $\mathcal{C} \subseteq \Sigma^\blen$ be a code, over any alphabet $\Sigma$, in which two distinct codewords agree on at most $D$ coordinates. For every integer $A$ with $A^2>\blen D$, every word $\mathbf{c} \in \Sigma^\blen$ has at most
\[
 \left\lfloor\frac{\blen(A-D)}{A^2-\blen D}\right\rfloor
\]
codewords agreeing with it on at least $A$ coordinates.
\end{theorem}

\begin{proof}
Let $U_1, \dots, U_L$ be distinct codewords, each agreeing with $\mathbf{c}$ on at least $A$ coordinates. For each $j$, fix a set $\aset_j$ of exactly $A$ coordinates on which $U_j$ agrees with $\mathbf{c}$, and write $\varphi:=A/\blen$ and $\rate:=D/\blen$. For a coordinate $i$, let $d_i := |\{ j : i \in \aset_j \}|$; then $\sum_i d_i = LA=L\varphi\blen$.

Wherever two codewords both agree with $\mathbf{c}$ they agree with each other, so $|\aset_j \cap \aset_{j'}| \le \rate \blen$ for $j \neq j'$, and counting ordered pairs,
\[
\sum_i d_i (d_i - 1) \;=\; \sum_{j \neq j'} |\aset_j \cap \aset_{j'}| \;\le\; L(L-1)\, \rate \blen .
\]
Cauchy--Schwarz on the vectors $(d_1, \dots, d_\blen)$ and $(1, \dots, 1)$ gives $\bigl(\sum_i d_i\bigr)^2 \le \blen \sum_i d_i^2$, i.e.\ $\sum_i d_i^2 \ge (L \varphi \blen)^2 / \blen = L^2 \varphi^2 \blen$; the left side above, which equals $\sum_i d_i^2 - \sum_i d_i$, is therefore at least $L^2 \varphi^2 \blen - L \varphi \blen$. Dividing by $L\blen$ and rearranging,
\[
L \varphi^2 - \varphi \;\le\; (L-1)\, \rate
\qquad\Longleftrightarrow\qquad
L\, (\varphi^2 - \rate) \;\le\; \varphi - \rate .
\]
Since $A^2>\blen D$, the denominator is positive, and the displayed inequality gives
\[
 L\le\frac{\varphi-\rate}{\varphi^2-\rate}=\frac{\blen(A-D)}{A^2-\blen D}.
\]
Taking the floor proves the claim.
\end{proof}
